\chapter{Variables, diagrams and justified disagreement}
\label{la:diagrams}

The preceding chapter established what definitions and fallible evidence can
contribute. Here those objects enter a reviewable diagram. The argument first
separates graph classes, premises and retrieval outcomes, then develops exact
proxy, collider, clock and disagreement examples. The ontology comparison
connects the diagram to implementable treatment versions and source records.
The final restrictions study follows a premise through both graph and effect
sets, showing why a smaller graph class need not settle the decision.

\section{From vocabulary to measured variables}

A useful causal diagram begins with variables whose meanings can survive
contact with data. A business glossary may contain ``risk,'' ``relationship''
and ``engagement,'' while a dataset contains a risk-score vintage, a set of
linked accounts and an observed click. Those are related concepts, but they
are not interchangeable causal variables. An assistant can propose mappings,
detect inconsistent names and assemble definitions. Each proposed mapping
still needs an operational meaning, a unit and an information clock.

OntoLearner offers relevant modular ideas for ontology loading, typing,
taxonomy and relation-learning tasks. The inspected paper reports
domain-dependent errors and retrieval failures. Its relation-learning
capability is not a guarantee that a relation is causal, or that a concept
has a valid measured proxy in the bank. Ontology alignment should therefore
precede and inform graph construction without being silently promoted into
causal identification. \citep{v12_la_2607_01977}

For example, ``customer owns account'' may describe a legal relation,
``records share an identifier'' an administrative relation, and ``records
have a high linkage score'' a model output. If ownership affects both offer
eligibility and repayment behavior, substituting the linkage score for
ownership leaves a measurement problem. Adding an arrow from the score to
repayment does not resolve it. The diagram should distinguish latent
ownership, observed records, the linkage procedure and sample selection.

Likewise, an ordinal description such as ``high influence'' cannot be used
as an effect size without a scale. An assistant's rating of a proposed edge
may express a preference among textual explanations. It is not automatically
a probability that the edge is present, a posterior over graphs, or a
confidence level. A typed diagram can retain such a rating as explicitly
uncalibrated metadata while refusing to use it as scientific evidence.

\section{What an edge asserts}

An edge proposal should state more than two variable names. The endpoint
marks identify the proposed graphical object. A source and inspected anchor
state what supports the proposal. A scope describes the population, period
and intervention to which that evidence applies. A counterevidence record
states whether contrary evidence was found, absent in the inspected material,
or not successfully searched. A review disposition records acceptance as a
premise, rejection, or unresolved status.

These fields serve different purposes. A source can be authentic but
inapplicable to the current banking action. A study can concern a different
population or a different definition of treatment. A plausible timing argument
can constrain direction without establishing that any causal connection
exists. The assistant should make the supported proposition as narrow as
the evidence warrants and retain the reasons for disagreement.

CausalChat supplies an example of interactive model development: the inspected
paper combines optional CPDAG discovery, model suggestions, graph editing
and structural-equation coefficient updates. Its prompts use ordinal
causal-significance ratings and ask for possible confounders and mediators.
These are useful interaction ideas. The paper's expert feedback also raises
missing citations and uncertainty about the quality of scientific evidence
provided by the model. A coefficient update after adding an edge cannot
establish that the edge itself is true. \citep{v12_la_2410_14146}

The proposed banking workflow retains the interaction but changes the
scientific handoff. An accepted premise is labeled as a premise with a
responsible review, not as a truth certified by the model. A model proposal
alone cannot remove an observationally compatible graph alternative. The
tool responsible for formal identification receives the complete reviewed
graph object and its restrictions, rather than a picture whose uncertainty
has disappeared in editing.

\section{Graph classes and collections of candidate models}

A CPDAG or PAG is a structured representation with its own interpretation.
A list of proposed DAGs is another object. A collection of edges marked
``possible'' is not automatically an enumeration of every combination of
those edges. Confusing these representations can silently exclude models
or include incompatible ones.

The project contract preserves a CPDAG or PAG directly when that is the
appropriate object, including circle endpoints. It also permits explicitly
listed candidate models. An edge marked possible within a candidate describes
the evidence status of that candidate's edge; it does not make all optional
edge combinations an exhaustive graph class. Downstream calculations must
know which interpretation applies.

An exact binary example shows the consequence of retaining alternatives.
In model \(F\), a binary offer \(A\) is Bernoulli with probability one half,
and response \(Y\) equals \(A\) with probability \(0.9\). In model \(R\),
\(Y\) is Bernoulli with probability one half, and \(A\) equals \(Y\) with
probability \(0.9\). Both models produce
\[
\begin{array}{c|cc}
 &Y=0&Y=1\\\hline
A=0&9/20&1/20\\
A=1&1/20&9/20
\end{array}.
\]
In \(F\), setting \(A\) from zero to one changes \(E[Y]\) by \(0.8\).
In \(R\), intervening on the equation for \(A\) leaves \(Y\) unchanged.
The two compatible constructions therefore imply the effect candidates
\(\{0,0.8\}\). A directional suggestion favoring \(F\) does not provide
new identifying information unless its underlying premise is independently
justified. These finite calculations are project constructions, not a
calibrated confidence set for a learned graph.

More generally, for a declared graph class \(\mathcal G\), one relevant
downstream object is
\begin{equation}
 \Theta(P,Q,\mathcal G)=\bigcup_{G\in\mathcal G}\Theta_G(P,Q),
 \label{la:graph-union}
\end{equation}
where \(\Theta_G(P,Q)\) contains effects compatible with distribution \(P\),
question \(Q\) and the model restrictions for \(G\). A nonidentified member
can make the union broad. A list of estimated points is not automatically
the whole union, and graph-selection uncertainty cannot be replaced with
ordinary coefficient uncertainty after selecting one member.

Equation~\eqref{la:graph-union} also clarifies what a disagreement review
should resolve. If two experts differ only on an edge that leaves the target
functional unchanged, that disagreement may be immaterial for the current
decision while still mattering for another target. If their graphs imply
different effect signs or feasibility, the conflict is decision-relevant.
The assistant can identify that consequence and prioritize the required
evidence. It should not average expert preferences into an unexplained graph.

\section{Retrieval failure and scientific disagreement}

An evidence assistant can fail to find a passage because the document is
unavailable, parsing failed, the query was poor, the retriever missed it, or
the passage is absent from the inspected document. These cases support
different statements. ``No passage was retrieved'' cannot by itself become
``the authors ignored the assumption,'' and neither statement establishes
that the assumption is false in the population.

The inspected ARGUS paper makes this distinction especially useful. It
audits reported evidence concerning DiD identifying assumptions and does
not claim to adjudicate causal truth. Its synthetic flaw study exhibits a
retrieval/abstention tradeoff, and its small human-label pilot reports an
in-sample calibration change that was not tested on unseen papers. These
observations motivate a banking evidence state that is separate from a
scientific risk judgment. \citep{v12_la_2609_30867}

The proposed record contains the source bytes, the query or retrieval
configuration, inspected anchors and the resulting narrow claim. If evidence
is missing, it records the missing step. If two sources conflict, it records
their populations, dates and definitions before treating the conflict as
substantive. The analyst can then determine whether the sources concern
different regimes or whether a genuine contradiction remains.

This process does not require pretending that every edge has a published
study behind it. Some causal premises depend on institutional facts: which
system sets a rate, when a policy becomes visible, how an audit was selected,
or which identifier establishes ownership. Those facts need an accountable
owner and authentic records. Literature retrieval can help formulate the
questions, but it cannot supply an institution's missing assignment truth.

\section{A reviewable editing sequence}

The proposed interaction starts with a question and variable definitions,
then records diagram alternatives. Each edit changes a declared object:
add or remove a variable, refine a clock, replace a proxy, propose an endpoint,
or accept a restriction. The edit retains its predecessor identity and
reason. A reviewer can therefore see whether an apparent improvement in
fit came from a new assumption, a different target, or an actual correction
of an error.

Formal checks can reject a directed cycle in a DAG, impossible declared
time order or an explicitly forbidden direction. They can flag mediator,
collider or common-ancestor patterns in the supported representation. Such
checks help maintain a coherent declaration but do not cover every graph
class or establish that an accepted edge is causally true. In particular,
a diagram with no syntax error can still omit an important common cause or
misrepresent the observation process.

The user should be able to retain an unresolved diagram and ask for analyses
conditional on several alternatives. A refusal to pick one graph can be the
scientifically correct output. The later discovery and identification
operations then need to report what changes across those alternatives and
what remains common. This makes uncertainty operational: it directs the
next information request rather than appearing only as a disclaimer after
a single confident answer.

The following constructions make that consequence explicit. They separate
observational equivalence from sampling precision, measured proxies from
confounders, and a graph edit from evidence that warrants it. Their monetary
comparisons then locate the disagreements worth investigating.
